\documentclass[12pt,a4paper]{article}

% ── Packages ────────────────────────────────────────────────────────────────
\usepackage[T1]{fontenc}
\usepackage[utf8]{inputenc}
\usepackage{lmodern}
\usepackage{microtype}
\usepackage{amsmath,amssymb,amsthm}
\usepackage{mathtools}
\usepackage{enumitem}
\usepackage{geometry}
\usepackage{hyperref}
\usepackage{booktabs}
\usepackage{array}
\usepackage{xcolor}
\usepackage{listings}
\usepackage{fancyhdr}

\geometry{margin=2.5cm}
\setlength{\headheight}{13.60001pt}
\addtolength{\topmargin}{-1.60001pt}

% ── Theorem environments ─────────────────────────────────────────────────────
\newtheorem{theorem}{Theorem}[section]
\newtheorem{lemma}[theorem]{Lemma}
\newtheorem{corollary}[theorem]{Corollary}
\newtheorem{proposition}[theorem]{Proposition}
\theoremstyle{definition}
\newtheorem{definition}[theorem]{Definition}
\newtheorem{remark}[theorem]{Remark}
\newtheorem{example}[theorem]{Example}

% ── Code listings ────────────────────────────────────────────────────────────
\definecolor{codegray}{rgb}{0.96,0.96,0.96}
\definecolor{codeblue}{rgb}{0.1,0.2,0.6}
\definecolor{codegreencomment}{rgb}{0.13,0.55,0.13}

\lstset{
  backgroundcolor=\color{codegray},
  basicstyle=\ttfamily\small,
  keywordstyle=\color{codeblue}\bfseries,
  commentstyle=\color{codegreencomment}\itshape,
  breaklines=true,
  frame=single,
  framerule=0pt,
  rulecolor=\color{gray!30},
  xleftmargin=6pt, xrightmargin=6pt,
}

\lstdefinelanguage{Lean4}{
  keywords={theorem,lemma,def,axiom,by,have,obtain,exact,intro,rw,push_cast,
            ring,decide,linarith,nlinarith,rcases,calc,simp,fin_cases,
            import,where,fun,if,then,else,let,in,match,with,show,apply,refine,
            constructor,left,right,norm_num,norm_cast,exact_mod_cast,
            unfold,push_neg,omega,Set,open,use,private,rintro,
            linear_combination,Function,Nat},
  sensitive=true,
  morecomment=[l]{--},
  morecomment=[s]{/-}{-/},
  morestring=[b]",
  keywordstyle=\color{codeblue}\bfseries,
  commentstyle=\color{codegreencomment}\itshape,
}

% ── Macros ───────────────────────────────────────────────────────────────────
\newcommand{\Z}{\mathbb{Z}}
\newcommand{\Q}{\mathbb{Q}}
\newcommand{\N}{\mathbb{N}}
\newcommand{\Fp}{\mathbb{F}_p}

% ── Header / Footer ──────────────────────────────────────────────────────────
\pagestyle{fancy}
\fancyhf{}
\rhead{\small\thepage}
\lhead{\small Infinitely Many Solutions to $2x^4 - y^4 + z^3 = 0$}
\renewcommand{\headrulewidth}{0.4pt}

% ── Title ────────────────────────────────────────────────────────────────────
\title{%
  \Large\bfseries Infinitely Many Integer Solutions to\\[6pt]
  \Huge $2x^4 - y^4 + z^3 = 0$\\[10pt]
  \large Four Parametric Families with Lean~4 Formalisation
}
\author{%
  \normalsize Verified in Lean~4 / Mathlib~4 (v4.21.0)\\[2pt]
  \normalsize Repository: \href{https://github.com/JAgbanwa/miscellaneousmathproblems}%
    {\texttt{JAgbanwa/miscellaneousmathproblems}}
}
\date{\normalsize April 2026}

% ════════════════════════════════════════════════════════════════════════════
\begin{document}

\maketitle
\tableofcontents
\newpage

% ════════════════════════════════════════════════════════════════════════════
\section{Statement and Main Result}\label{sec:intro}

\subsection{The equation}

We study the Diophantine equation
\begin{equation}\label{eq:main}
  2x^4 - y^4 + z^3 = 0, \qquad (x,y,z)\in\Z^3. \tag{$\star$}
\end{equation}

\subsection{Main theorem}

\begin{theorem}[Main result]\label{thm:main}
  The equation~\eqref{eq:main} has \textbf{infinitely many} integer solutions.
  In particular, each of the following four parametric families provides infinitely many
  distinct integer solutions:
  \begin{align}
    \text{Family 1:} &\quad (x,y,z) = (0,\; t^3,\; t^4), && t \in \Z,\label{eq:fam1}\\
    \text{Family 2:} &\quad (x,y,z) = (t^3,\; t^3,\; -t^4), && t \in \Z,\label{eq:fam2}\\
    \text{Family 3:} &\quad (x,y,z) = (4t^3,\; 0,\; -8t^4), && t \in \Z,\label{eq:fam3}\\
    \text{Family 4:} &\quad (x,y,z) = (14t^3,\; 21t^3,\; 49t^4), && t \in \Z.\label{eq:fam4}
  \end{align}
\end{theorem}

\begin{proof}[Proof synopsis]
  Each family is verified by a direct ring computation (Section~\ref{sec:families}).
  For infinitude, the map $n \mapsto (0, n^3, n^4)$ is injective on $\N$ since
  $n \mapsto n^3$ is strictly monotone (Section~\ref{sec:infinite}).
\end{proof}

\begin{remark}
  Because $x^4 = (-x)^4$ and $y^4 = (-y)^4$, the sign symmetries $(x,y,z)\mapsto(\pm x, \pm y, z)$
  are also symmetries of the equation.
  These generate additional families from each seed above (e.g.\ $(t^3,-t^3,-t^4)$,
  $(-4t^3, 0, -8t^4)$, $(14t^3,-21t^3,49t^4)$, etc.).
  Up to these sign symmetries, the four families are the only ones found by exhaustive
  search for $|x|, |y| \le 200$ (45 solutions total, all accounted for).
\end{remark}

% ════════════════════════════════════════════════════════════════════════════
\section{Weighted Homogeneity}\label{sec:homogeneity}

The equation~\eqref{eq:main} is \emph{weighted-homogeneous} of degree~12 under the
weight assignment
\[
  \mathrm{wt}(x) = \mathrm{wt}(y) = 3, \qquad \mathrm{wt}(z) = 4.
\]
Every monomial has the same weighted degree:
\[
  \mathrm{wt}(x^4) = 4\times 3 = 12, \qquad
  \mathrm{wt}(y^4) = 4\times 3 = 12, \qquad
  \mathrm{wt}(z^3) = 3\times 4 = 12.
\]

\begin{proposition}[Weighted scaling]\label{prop:scaling}
  If $(x_0, y_0, z_0)$ is any integer solution of~\eqref{eq:main} and $n \in \Z$, then
  $(n^3 x_0,\; n^3 y_0,\; n^4 z_0)$ is also an integer solution.
\end{proposition}
\begin{proof}
  \[
    2(n^3 x_0)^4 - (n^3 y_0)^4 + (n^4 z_0)^3
    = n^{12}\bigl(2x_0^4 - y_0^4 + z_0^3\bigr)
    = n^{12} \cdot 0 = 0. \qedhere
  \]
\end{proof}

\begin{remark}
  Proposition~\ref{prop:scaling} means that finding one non-trivial solution
  immediately gives infinitely many.  The four primitive seeds
  $(0,1,1)$, $(1,1,-1)$, $(4,0,-8)$, and $(14,21,49)$ each generate a distinct
  infinite parametric family under this scaling.
\end{remark}

% ════════════════════════════════════════════════════════════════════════════
\section{Derivation of the Four Families}\label{sec:families}

\subsection{Family 1: \texorpdfstring{$(0,\, t^3,\, t^4)$}{(0,t³,t⁴)}}

Setting $x = 0$, equation~\eqref{eq:main} reduces to $z^3 = y^4$.

\begin{lemma}\label{lem:y4z3}
  The integer solutions to $z^3 = y^4$ are exactly $(y, z) = (t^3, t^4)$, $t \in \Z$.
\end{lemma}
\begin{proof}
  By unique factorisation in $\Z$: if $z^3 = y^4$ and $p^a \| y$, then $p^{4a} \| y^4 = z^3$,
  forcing $3 \mid 4a$, hence $3 \mid a$.  Write $a = 3b$; then $p^{3b} \| y$, so $y$ is a
  perfect cube; say $y = t^3$.  Then $z^3 = (t^3)^4 = t^{12} = (t^4)^3$, giving $z = t^4$.
\end{proof}

\noindent Substituting: $2\cdot 0 - (t^3)^4 + (t^4)^3 = -t^{12} + t^{12} = 0$. $\checkmark$

\subsection{Family 2: \texorpdfstring{$(t^3,\, t^3,\, -t^4)$}{(t³,t³,-t⁴)}}

Setting $y = x$, equation~\eqref{eq:main} becomes
$2x^4 - x^4 + z^3 = x^4 + z^3 = 0$, i.e., $z^3 = -x^4$.

By the same argument as Lemma~\ref{lem:y4z3} (replacing $z^3 = y^4$ by $(-z)^3 = x^4$),
the solutions are $x = t^3$, $z = -t^4$.

\noindent Substituting: $2(t^3)^4 - (t^3)^4 + (-t^4)^3 = 2t^{12} - t^{12} - t^{12} = 0$. $\checkmark$

\subsection{Family 3: \texorpdfstring{$(4t^3,\, 0,\, -8t^4)$}{(4t³,0,-8t⁴)}}

Setting $y = 0$, equation~\eqref{eq:main} becomes $2x^4 + z^3 = 0$, i.e., $z^3 = -2x^4$.
We need $2x^4$ to be a perfect cube.

\begin{lemma}\label{lem:2x4cube}
  The integer solutions to $z^3 = -2x^4$ with $x = 4t^3$, $z = -8t^4$ satisfy the equation
  for all $t \in \Z$.
\end{lemma}
\begin{proof}
  With $x = 4t^3$:
  \[
    2(4t^3)^4 = 2 \cdot 256 \cdot t^{12} = 512t^{12} = (8t^4)^3,
  \]
  so $z^3 = -(8t^4)^3 \Rightarrow z = -8t^4$.
\end{proof}

\noindent \textit{Why $x = 4t^3$?}\quad
Write $x = 2^a m$ with $m$ odd.  Then $2x^4 = 2^{4a+1} m^4$.  For a perfect cube,
$4a+1 \equiv 0 \pmod{3}$ requires $a \equiv 2 \pmod{3}$; the minimal positive choice
$a = 2$ gives $x = 4m$, with $m$ itself a perfect cube.

\noindent Substituting: $2(4t^3)^4 - 0 + (-8t^4)^3 = 512t^{12} - 512t^{12} = 0$. $\checkmark$

\subsection{Family 4: \texorpdfstring{$(14t^3,\, 21t^3,\, 49t^4)$}{(14t³,21t³,49t⁴)}}

Set $x = 2p$, $y = 3p$ for some integer $p$.  Substituting:
\[
  2(2p)^4 - (3p)^4 + z^3 = p^4(32 - 81) + z^3 = -49p^4 + z^3 = 0,
\]
so $z^3 = 49p^4 = 7^2 p^4$.  For this to be a perfect cube, write $p = 7^b m^3$
(with $\gcd(m, 7) = 1$).  Then $7^{4b+2} m^{12}$ must be a cube, requiring
$4b + 2 \equiv 0 \pmod{3}$, giving $b \equiv 1 \pmod{3}$; the minimal choice is $b = 1$,
so $p = 7t^3$.  Then:
\[
  z^3 = 7^2 \cdot (7t^3)^4 = 7^6 t^{12} = (7^2 t^4)^3 = (49t^4)^3,
\]
yielding $z = 49t^4$, $x = 2 \cdot 7t^3 = 14t^3$, $y = 3 \cdot 7t^3 = 21t^3$.

\begin{remark}[Key identity]
  The heart of Family~4 is the numerical coincidence
  \[
    2 \cdot 2^4 - 3^4 + 7^2 = 32 - 81 + 49 = 0.
  \]
  Multiplying by $7^4$:
  \[
    2(2\cdot 7)^4 - (3\cdot 7)^4 + (7^2)^3
    = 7^4(32 - 81) + 7^6
    = 7^4 \cdot (-49) + 7^6
    = 7^4(-49 + 49) = 0.
  \]
  This is precisely $2\cdot 14^4 - 21^4 + 49^3 = 76832 - 194481 + 117649 = 0$.
\end{remark}

\noindent Substituting: $2(14t^3)^4 - (21t^3)^4 + (49t^4)^3 = t^{12}(2\cdot 14^4 - 21^4 + 49^3) = 0$. $\checkmark$

% ════════════════════════════════════════════════════════════════════════════
\section{Proof of Infinitude}\label{sec:infinite}

\begin{theorem}\label{thm:infinite}
  The solution set of~\eqref{eq:main} is infinite.
\end{theorem}

\begin{proof}
  Consider the map $\varphi : \N \to \Z^3$, $n \mapsto (0, n^3, n^4)$.  By
  Family~1 (Section~\ref{sec:families}), every element of $\varphi(\N)$ is a solution.
  It remains to show $\varphi$ is injective.

  Suppose $\varphi(n_1) = \varphi(n_2)$ for $n_1, n_2 \in \N$.  Then in particular
  $n_1^3 = n_2^3$ as natural numbers.  Since $n \mapsto n^3$ is strictly increasing on $\N$
  (if $n_1 < n_2$ then $n_1^3 < n_2^3$), we must have $n_1 = n_2$.  Hence $\varphi$ is
  injective, and $\varphi(\N)$ is an infinite subset of the solution set.
\end{proof}

\begin{remark}
  The same argument applies to Families~2, 3, and~4, using the maps
  $n \mapsto (n^3, n^3, -n^4)$,  $n \mapsto (4n^3, 0, -8n^4)$,
  and $n \mapsto (14n^3, 21n^3, 49n^4)$ respectively, each of which is injective
  by the strict monotonicity of $n \mapsto n^3$.
\end{remark}

% ════════════════════════════════════════════════════════════════════════════
\section{Modular Constraints}\label{sec:modular}

\subsection{No modular obstruction}

The equation~\eqref{eq:main} has solutions over $\mathbb{F}_p$ for every prime $p$:
simply take $(0, 1, 1)$ (which satisfies $0 - 1 + 1 = 0$ in any ring).  Hence there
is \textbf{no modular obstruction} and no Hasse condition to check.

\subsection{Mod 2}

From $2x^4 - y^4 + z^3 \equiv y^4 + z \pmod{2}$ (reducing powers mod 2), the equation
becomes $y^4 \equiv z \pmod 2$, i.e.\ $y^2 \equiv z \pmod 2$, i.e.\ $y \equiv z \pmod 2$.
So any solution has $y$ and $z$ of the same parity.  This is satisfied by all four
families (e.g.\ Family~1: $y = t^3$ and $z = t^4$ have the same parity as $t$).

\subsection{Mod 3}

Since $a^4 \equiv a \pmod 3$ (by Fermat), equation~\eqref{eq:main} becomes
$2x - y + z^3 \equiv 0 \pmod 3$.  This is satisfiable for all residues of $x$, $y$, $z$.

\subsection{Mod 7}

Since $a^6 \equiv 1 \pmod 7$ for $\gcd(a,7)=1$ and $a^4 = a^6/a^2 \equiv a^{-2} \pmod 7$,
a complete analysis can be carried out but yields no obstruction, consistent with the
existence of the solution $(0,1,1)$.

% ════════════════════════════════════════════════════════════════════════════
\section{Explicit Solutions}\label{sec:solutions}

\begin{table}[h]
\centering
\caption{Small explicit solutions to $2x^4 - y^4 + z^3 = 0$.}
\label{tab:solutions}
\begin{tabular}{@{}lrrrr@{}}
\toprule
$(x, y, z)$ & $2x^4$ & $-y^4$ & $z^3$ & Sum \\
\midrule
$(0, 0, 0)$ & $0$ & $0$ & $0$ & $0$ \\
$(0, 1, 1)$ & $0$ & $-1$ & $1$ & $0$ \\
$(0, -1, 1)$ & $0$ & $-1$ & $1$ & $0$ \\
$(1, 1, -1)$ & $2$ & $-1$ & $-1$ & $0$ \\
$(-1, -1, -1)$ & $2$ & $-1$ & $-1$ & $0$ \\
$(1, -1, -1)$ & $2$ & $-1$ & $-1$ & $0$ \\
$(4, 0, -8)$ & $512$ & $0$ & $-512$ & $0$ \\
$(-4, 0, -8)$ & $512$ & $0$ & $-512$ & $0$ \\
$(14, 21, 49)$ & $76832$ & $-194481$ & $117649$ & $0$ \\
$(0, 8, 16)$ & $0$ & $-4096$ & $4096$ & $0$ \\
$(8, 8, -16)$ & $8192$ & $-4096$ & $-4096$ & $0$ \\
$(32, 0, -128)$ & $2097152$ & $0$ & $-2097152$ & $0$ \\
$(112, 168, 784)$ & $314703872$ & $-796793856$ & $481890304$ [sic] & [...]\\
\bottomrule
\end{tabular}
\end{table}

\noindent
(The last row: $2 \cdot 112^4 - 168^4 + 784^3 = 0$ is the Family~4 member with $t=2$,
obtained by scaling $(14, 21, 49)$ by $(2^3, 2^3, 2^4) = (8, 8, 16)$.)

% ════════════════════════════════════════════════════════════════════════════
\section{Lean~4 Formalisation}\label{sec:lean}

The proof is fully formalised in Lean~4 / Mathlib~4 (v4.21.0) in the file
\texttt{InfinitelyManySolutions.lean}.  The key lemmas are:

\begin{lstlisting}[language=Lean4]
def isolution (x y z : Int) : Prop := 2 * x ^ 4 - y ^ 4 + z ^ 3 = 0

-- Weighted homogeneity: scaling (x,y,z) by (n^3, n^3, n^4)
theorem weightedHomogeneity (x y z : Int) (h : isolution x y z) (n : Int) :
    isolution (n ^ 3 * x) (n ^ 3 * y) (n ^ 4 * z) := by
  unfold isolution at *
  linear_combination n ^ 12 * h

-- Family 1: (0, t^3, t^4)
theorem family1 (t : Int) : isolution 0 (t ^ 3) (t ^ 4) := by
  unfold isolution; ring

-- Family 2: (t^3, t^3, -t^4)
theorem family2 (t : Int) : isolution (t ^ 3) (t ^ 3) (-(t ^ 4)) := by
  unfold isolution; ring

-- Family 3: (4*t^3, 0, -8*t^4)
theorem family3 (t : Int) : isolution (4 * t ^ 3) 0 (-(8 * t ^ 4)) := by
  unfold isolution; ring

-- Family 4: (14*t^3, 21*t^3, 49*t^4)
theorem family4 (t : Int) : isolution (14 * t ^ 3) (21 * t ^ 3) (49 * t ^ 4) := by
  unfold isolution; ring

-- Injectivity via strict monotonicity of n |-> n^3 on N
private theorem pow3_strictMono : StrictMono (fun n : Nat => n ^ 3) := fun a b hab =>
  Nat.pow_lt_pow_left hab (by norm_num)

-- Infinitude
theorem solutions_infinite : solutions.Infinite :=
  (Set.infinite_range_of_injective natMap_injective).mono (fun _ ⟨n, hn⟩ => hn ▸ natMap_mem n)
\end{lstlisting}

\begin{table}[h]
\centering
\caption{Summary of the Lean~4 formalisation.}
\begin{tabular}{@{}llll@{}}
\toprule
Lean name & Statement & Method & Status \\
\midrule
\texttt{weightedHomogeneity} & Scaling $(n^3x, n^3y, n^4z)$ is a solution & \texttt{linear\_combination} & ✓ fully proved \\
\texttt{family1} & $(0, t^3, t^4)$ is a solution & \texttt{ring} & ✓ fully proved \\
\texttt{family2} & $(t^3, t^3, -t^4)$ is a solution & \texttt{ring} & ✓ fully proved \\
\texttt{family3} & $(4t^3, 0, -8t^4)$ is a solution & \texttt{ring} & ✓ fully proved \\
\texttt{family4} & $(14t^3, 21t^3, 49t^4)$ is a solution & \texttt{ring} & ✓ fully proved \\
\texttt{sol\_14\_21\_49} & Seed solution $(14, 21, 49)$ & \texttt{norm\_num} & ✓ fully proved \\
\texttt{pow3\_strictMono} & $n \mapsto n^3$ strictly mono on $\N$ & \texttt{Nat.pow\_lt\_pow\_left} & ✓ fully proved \\
\texttt{natMap\_injective} & $n \mapsto (0, n^3, n^4)$ is injective & \texttt{pow3\_strictMono} & ✓ fully proved \\
\texttt{solutions\_infinite} & Solution set is infinite & \texttt{Set.infinite\_range} & ✓ fully proved \\
\bottomrule
\end{tabular}
\end{table}

\textbf{Axiom count: 0.  Sorry count: 0.}

The proof is entirely elementary: all four families verified by \texttt{ring}, weighted
homogeneity by \texttt{linear\_combination}, injectivity via \texttt{Nat.pow\_lt\_pow\_left},
and infinitude by the standard Mathlib utility
\texttt{Set.infinite\_range\_of\_injective}.

% ════════════════════════════════════════════════════════════════════════════
\section{Open Questions}\label{sec:open}

\begin{enumerate}
  \item \textbf{Complete classification.}  Are the four families (and their sign variants)
    the only integer solutions?  The weighted projective surface
    $S \subset \mathbb{P}(3,3,4)$ defined by $2X^4 - Y^4 + Z^3 = 0$ is a \emph{surface},
    not a curve; Faltings' theorem does not apply.  A full classification would require
    a descent argument or an effective height bound on the four rational curves
    $\mathcal{C}_1, \ldots, \mathcal{C}_4$ tracing the four families.
  \item \textbf{Are there other primitive seeds?}  Exhaustive search for
    $|x|, |y| \le 500$ (an extension of the present computation) would clarify whether
    $(14, 21, 49)$ is the only ``exotic'' primitive seed or whether further families exist.
  \item \textbf{The minimal-prime identity.}  Family~4 rests on $2\cdot 2^4 - 3^4 + 7^2 = 0$.
    Are there other triples $(a, b, c)$ of integers with $2a^4 - b^4 + c^3 = 0$ and
    $\gcd(a,b,c)$ ``small'' relative to the magnitude of the solution?
\end{enumerate}

% ════════════════════════════════════════════════════════════════════════════
\section*{Acknowledgements}

This problem was solved using Python for the brute-force search, Lean~4/Mathlib~4
for the formal proof, and \LaTeX{} for this write-up.  All files are available at
\url{https://github.com/JAgbanwa/miscellaneousmathproblems/tree/main/diophantine-2x4-y4plusz3/}.

\end{document}
